%% kapa: exam=1 round=20 year=2009 period=1 form=A subject=civil_law subject_ko=민법 seq=22 number=22
%% source: Q-Net 감정평가사 기출문제 (cst00309) artlSeq=5024178 file=extracted/1008/5024178/2009년 제20회 감정평가사 1차시험 기출문제/1교시.hwp
%% conversion: kordoc 4.24.51 -> tools/gen1.py (auto); review=proofread-sonnet; flags=none
\begin{question}{22}
집합건물의 소유 및 관리에 관한 설명으로 옳지 \uline{않은} 것은? (다툼이 있으면 판례에 의함)
\begin{choices}
  \item 전유부분이 수인의 공유에 속하는 경우에 공유자는 관리단집회에서 각각 의결권을 행사할 수 있다.
  \item 재건축 비용의 분담에 관한 사항을 정하지 아니한 재건축 결의는 특별한 사정이 없는 한 무효이다.
  \item 대지(垈地) 위에 구분소유권의 목적인 건물이 속하는 1동의 건물이 있을 때에는 그 대지의 공유자는 그 건물의 사용에 필요한 범위 내의 대지에 대하여 분할을 청구하지 못한다.
  \item 각 공유자는 규약에 달리 정함이 없는 한 그 지분의 비율에 따라 공용부분의 관리비용 기타 의무를 부담한다.
  \item 집합건물구분소유권의 특별승계인이 구분소유권을 다시 제3자에게 이전한 경우에도 자신의 전(前)구분소유자의 공용부분에 대한 체납관리비를 지급할 책임이 있다.
\end{choices}
\end{question}
\answer{1}
